\chapter{Medicines for bipolar disorder}
\index{Medicines for bipolar disorder}

\section*{Definition and Overview}
Medicines for bipolar disorder refer to a diverse group of psychiatric medications used to stabilise mood fluctuations, prevent relapses, and manage acute episodes of mania, hypomania, and depression in patients diagnosed with bipolar I or bipolar II disorder. The primary classes of medications include mood stabilisers (such as lithium and sodium valproate), atypical antipsychotics (such as quetiapine, olanzapine, and aripiprazole), and certain anticonvulsants (such as lamotrigine). These drugs target neurotransmitter systems in the brain to restore neurochemical balance and enhance neuroplasticity. Because bipolar disorder is a lifelong, cyclic psychiatric condition, pharmacological treatment is typically long-term and requires regular clinical monitoring, blood tests, and adjustments to ensure therapeutic efficacy and minimise metabolic, renal, or endocrine side effects.

\section*{Epidemiology}
Bipolar disorder affects approximately 1\% to 2\% of the global population, with bipolar I and II presenting equal prevalence rates. In many countries, it is estimated that around 1 in 50 individuals will experience bipolar disorder in their lifetime. Pharmacological management is the cornerstone of treatment for almost all diagnosed patients. The peak onset of bipolar disorder is in late adolescence or early adulthood (ages 15 to 25). The lifetime adherence rate for bipolar medications is a significant clinical challenge, with up to 50\% of patients discontinuing therapy at some point due to side effects, lack of insight during manic phases, or the burden of long-term monitoring, leading to a high rate of recurrence and hospitalisation.

\section*{Aetiology and Pathophysiology}
Bipolar disorder is associated with complex neurobiological alterations, and medications work by modulating these disrupted pathways:
\begin{itemize}
    \item \textbf{Neurotransmitter dysregulation:} Abnormalities in monoaminergic systems (dopamine, serotonin, noradrenaline). Antipsychotics act as antagonists at dopamine D2 and serotonin 5-HT2A receptors to control manic symptoms.
    \item \textbf{Intracellular signalling pathways:} Lithium and sodium valproate inhibit glycogen synthase kinase-3 (GSK-3) and inositol monophosphatase, modulating intracellular signalling cascades and exerting neuroprotective effects.
    \item \textbf{Ion channel stabilisation:} Anticonvulsant mood stabilisers (like lamotrigine or carbamazepine) block voltage-gated sodium channels, reducing the release of excitatory neurotransmitters like glutamate.
    \item \textbf{GABAergic neurotransmission:} Sodium valproate increases gamma-aminobutyric acid (GABA) levels, reinforcing inhibitory pathways to suppress neuronal hyperexcitability during manic episodes.
    \item \textbf{HPG and HPA axis dysfunction:} Neuroendocrine dysregulation is modulated indirectly through long-term mood stabilisation.
\end{itemize}

\section*{Clinical Features}
Patients requiring or undergoing bipolar medication monitoring display clinical presentations related to their mood state or medication-induced side effects:
\begin{itemize}
    \item \textbf{Acute manic features:} Pressured speech, grandiosity, decreased need for sleep, racing thoughts, and risky behaviours, requiring acute treatment with antipsychotics or lithium.
    \item \textbf{Acute depressive features:} Profound sadness, anhedonia, fatigue, and suicidal ideation, which may be managed with lamotrigine or quetiapine (avoiding antidepressant monotherapy to prevent switching to mania).
    \item \textbf{Euthymic (stable) state:} The target therapeutic state, presenting with normal social and occupational functioning and absence of distinct mood episodes.
    \item \textbf{Common medication side effects:} Weight gain, sedation, hand tremors, polyuria, polydipsia, and gastrointestinal disturbances (nausea or diarrhoea).
    \item \textbf{Signs of toxicity (especially lithium):} Coarse tremor, ataxia, slurred speech, confusion, vomiting, and severe muscle weakness.
\end{itemize}

\section*{Differential Diagnosis}
Side effects or symptoms during treatment must be distinguished from other medical and psychiatric conditions:
\begin{itemize}
    \item \textbf{Lithium toxicity vs. Gastrointestinal infection:} Nausea, vomiting, and diarrhoea can indicate acute lithium toxicity, which must be verified with urgent serum lithium levels.
    \item \textbf{Medication-induced sedation vs. Depressive relapse:} Sleepiness and lack of energy from high doses of antipsychotics can mimic depressive symptoms.
    \item \textbf{Neuroleptic malignant syndrome (NMS) vs. Serotonin syndrome:} Rare, life-threatening reactions to antipsychotics or combined drug therapies, presenting with rigidity, hyperthermia, and autonomic instability.
    \item \textbf{Drug-induced tremor vs. Essential tremor:} Fine postural tremors induced by lithium or valproate, which may respond to dose reduction or beta-blockers.
    \item \textbf{Thyroid dysfunction vs. Mood symptoms:} Lithium-induced hypothyroidism presenting with fatigue and weight gain, mimicking depression.
\end{itemize}

\section*{When to Seek Medical Care}
Patients taking bipolar medicines should have regular medical reviews. Family members and patients must be aware of critical emergency symptoms.

\textbf{Contact your local emergency services for:}
\begin{itemize}
    \item \textbf{Acute lithium overdose or toxicity:} Severe confusion, slurred speech, inability to walk (ataxia), vomiting, or seizures.
    \item \textbf{Signs of NMS:} High fever, severe muscle stiffness ("lead-pipe" rigidity), rapid heart rate, and confusion after starting or increasing antipsychotics.
    \item \textbf{Severe cutaneous reactions:} A spreading, painful rash with blisters and mucosal involvement (suspected Stevens-Johnson syndrome, a rare risk associated with lamotrigine).
    \item \textbf{Acute suicidal crisis:} Clear verbalisations or plans of self-harm or suicide.
\end{itemize}
Other reasons to contact the doctor include a fine tremor that becomes a coarse shaking, significant weight gain or increased thirst, or plans to become pregnant (as several bipolar medications, especially sodium valproate, carry high risks of birth defects).

\section*{Diagnosis and Investigations}
Safe prescribing of bipolar medicines requires comprehensive baseline and ongoing monitoring:
\begin{itemize}
    \item \textbf{Serum lithium levels:} Checked regularly (every 3 to 6 months) via a blood test taken exactly 12 hours after the last dose, with a narrow target therapeutic range of 0.6 to 0.8 mmol/L (up to 1.0 mmol/L in acute mania).
    \item \textbf{Renal function tests (eGFR and creatinine):} Performed before starting lithium and every 3 to 6 months to monitor for lithium-induced chronic kidney disease.
    \item \textbf{Thyroid function tests (TSH and free T4):} Monitored every 6 to 12 months, as lithium can cause goitre and hypothyroidism.
    \item \textbf{Liver function tests (LFTs) and Full blood count:} Indicated for patients taking sodium valproate to monitor for hepatotoxicity and thrombocytopenia.
    \item \textbf{Metabolic monitoring:} Regular assessment of weight, waist circumference, fasting lipids, and HbA1c, particularly for patients on atypical antipsychotics due to risk of metabolic syndrome.
    \item \textbf{Pregnancy test:} Mandatory for female patients of childbearing potential before initiating sodium valproate or lithium.
\end{itemize}

\section*{Management}
Pharmacological treatment must be individualised and integrated with psychoeducation:
\begin{itemize}
    \item \textbf{Mood stabilisers (lithium):} The gold standard for relapse prevention and managing mania. Sodium valproate is an alternative, especially for rapid-cycling bipolar disorder.
    \item \textbf{Atypical antipsychotics:} Used for acute mania (e.g. aripiprazole, risperidone) and bipolar depression (e.g. quetiapine, lurasidone).
    \item \textbf{Anticonvulsant mood stabilisers (lamotrigine):} Highly effective for preventing depressive episodes, titrated slowly over several weeks to minimise rash risk.
    \item \textbf{Psychoeducation and therapy:} Structured programmes to help patients understand their illness, identify early warning signs of relapse (prodromes), and improve medication adherence.
\end{itemize}

\section*{Complications}
Potential complications associated with long-term bipolar pharmacotherapy include:
\begin{itemize}
    \item \textbf{Nephrogenic diabetes insipidus:} Lithium-induced resistance to ADH, causing excessive urination (polyuria) and thirst.
    \item \textbf{Chronic kidney disease (CKD):} Progressive decline in glomerular filtration rate after decades of lithium therapy.
    \item \textbf{Metabolic syndrome and type 2 diabetes:} Secondary to antipsychotic-induced weight gain and insulin resistance.
    \item \textbf{Teratogenicity:} Fetal valproate syndrome (neural tube defects, cognitive impairment) and Ebstein's anomaly (a rare cardiac defect associated with lithium use in early pregnancy).
    \item \textbf{Tardive dyskinesia:} Involuntary, repetitive facial and limb movements secondary to long-term antipsychotic therapy.
\end{itemize}

\section*{Prognosis and Outcomes}
The prognosis for individuals with bipolar disorder is significantly improved by consistent, long-term adherence to appropriate medication. Pharmacotherapy can reduce the frequency and severity of mood episodes by up to 60\% to 80\%, allowing many patients to lead stable, productive lives. Finding the right medication combination and dosage can take time, and close collaboration between the patient, their psychiatrist, and general practitioner is essential. When combined with supportive psychotherapy, lifestyle stability (e.g. regular sleep patterns), and avoidance of substance abuse, the long-term prognosis for maintaining euthymia is high.

\section*{Prevention and Self-Care}
Adherence to self-care strategies is critical for preventing relapses and managing medication side effects:
\begin{itemize}
    \item \textbf{Maintain a strict dosing schedule:} Take medications at the same time every day, and never stop taking them suddenly without consulting your doctor.
    \item \textbf{Stay hydrated:} Drink a consistent amount of water daily, especially when taking lithium, and avoid sudden changes in salt intake, as dehydration or low-salt diets can cause toxic lithium accumulation.
    \item \textbf{Avoid alcohol and recreational drugs:} These substances can trigger mood episodes and interact negatively with bipolar medications.
    \item \textbf{Track mood and sleep patterns:} Use a mood diary or app to log daily mood, sleep duration, and medication compliance, helping to detect early signs of mania or depression.
    \item \textbf{Attend all laboratory monitoring sessions:} Complete scheduled blood tests for lithium levels, kidney, and thyroid function as requested.
\end{itemize}

\section*{Sources and Further Reading}
The following reputable organisations provide further information on bipolar medications:
\begin{itemize}
    \item Black Dog Institute. \textit{Bipolar disorder treatments and medications.} https://www.blackdoginstitute.org.au/
    \item SANE Australia. \textit{Bipolar disorder guide and support.} https://www.sane.org/
    \item Better Health Channel. \textit{Bipolar disorder medical treatments.} https://www.betterhealth.vic.gov.au/
    \item Mayo Clinic. \textit{Bipolar disorder: diagnosis, drugs, and therapy.} https://www.mayoclinic.org/
    \item NHS. \textit{Treatment of bipolar disorder: medications and monitoring.} https://www.nhs.uk/
\end{itemize}
